We evaluate the composed system along two complementary axes: controlled fault injection, which provides per-type ground truth, and a real fault deployment, which provides an external validity check. Because no public stream simultaneously labels all four fault types per signal and per timestep, neither axis alone suffices, and we triangulate the two. All evaluation artifacts referenced below are committed as CSV files under \texttt{examples/benchmarks/}.

\subsection{Controlled fault injection}
Following the injection-protocol precedent of~\citet{Lai2021}, we inject parameterized faults of each type---bias, drift, accuracy loss, and freezing---into a multivariate process with known ground truth, drawn from the Controlled Anomalies Time Series (CATS) benchmark~\citep{CATS2023}. Each typed fault is injected at controlled magnitude and onset onto a target signal, and recall is measured per type. We run $N=85$ trials per type in the scaled controlled validation, and report Wilson 95\% confidence intervals on every recall and precision estimate (\texttt{i7\_scaled\_per\_type.csv}). To probe the dependence of bias recall on conditional coupling, we report bias recall per target channel against that channel's conditional standard deviation (\texttt{i7\_scaled\_bias\_by\_channel.csv}, $N=17$ per channel). Freeze attribution versus base detection is evaluated on the same channels, recording per-channel freeze and CDF detection rates and detection delays (\texttt{i7\_freeze\_latency.csv}).

\subsection{Synthetic correlation control}
To separate the two robustness mechanisms of the composed system, we generate synthetic multivariate streams with controlled pairwise correlation $\rho \in \{0.1, 0.3, 0.5, 0.7, 0.9\}$ and inject coordinated shifts of magnitude $\{3, 6, 10\}$ conditional standard deviations co-occurring with a change point. For each cell we sweep the suppression scale $\texttt{suppress\_scale} \in \{1, 5, \infty\}$ (abbreviated \texttt{scale} hereafter) and report the mean false-label rate over $5$ seeds (\texttt{i7\_regime\_fp.csv}). The co-occurring-fault dead-band is measured by injecting a drift simultaneously with a change point and recording the detection magnitude at each suppression scale (\texttt{i7\_deadband.csv}). The freeze-on-quiescent ablation injects no fault into a genuinely live but quiet signal (noise floor as a fraction of marginal standard deviation) and measures the false freeze-label rate under default, strict, and loose freeze constants (\texttt{i7\_freeze\_quiescent.csv}).

\subsection{Real fault deployment}
For external validity we use the Intel Berkeley Lab sensor deployment~\citep{IntelLab2004}, in which motes report temperature and humidity over a diurnal cycle. As a mote's battery depletes (voltage below $2.3$\,V, held out as a physical fault label), its temperature and humidity readings degrade to physically impossible values---temperatures of $100$--$122$\,$^{\circ}$C and negative humidity. We treat the battery-voltage label as ground truth for a real fault, never exposing it to the detector or classifier. We sweep the detector's mean-threshold operating point and report, pooled over motes, the healthy false-positive rate and the real-fault detection rate ($N_\mathrm{healthy} \approx 1.52$M, $N_\mathrm{fault} \approx 0.39$M; \texttt{i7\_intel\_lab\_operating\_points.csv}), together with the distribution of types assigned to detected real faults (\texttt{i7\_intel\_lab\_fault\_types.csv}). To establish that online adaptation is required rather than optional, we additionally fit a static model on a short prefix of the same stream and measure its false-alarm rate.

As a second, complementary real-fault check that isolates a single labelled type, we use the LBNL building-FDD MZVAV-1 dataset~\citep{LBNL_FDD2019,Granderson2023}, an outdoor-air-temperature sensor-bias set in which a constant $\pm 1/2/4\,^\circ$C offset is imposed on the outdoor-air-temperature reading of a simulated air-handling unit over labelled date windows, with a per-timestamp fault label. We run the detector over the four coupled AHU air temperatures and report, over a mean-threshold sweep, the bias-typing recall on the biased signal against the healthy false-positive rate, together with the imposed bias magnitude in conditional-standard-deviation units (\texttt{i7\_lbnl\_bias.csv}).

\subsection{Metrics}
For the controlled injections we report per-type recall (fraction of injected faults of a type that are detected) and cross-talk precision on normal peer signals. For the real deployment we report the healthy false-positive rate and the real-fault detection rate at each operating point. All proportions estimated from finite samples are reported with Wilson 95\% confidence intervals. We do not report aggregate accuracy across types, because the type populations are not equiprobable in either the injection design or the real stream, and an aggregate would obscure the per-type and per-coupling structure that is the substance of our characterization.
